\documentclass[11pt, a4paper]{article}
\usepackage[utf8]{inputenc}
\usepackage[T1]{fontenc}
\usepackage[margin=1in]{geometry}
\usepackage{graphicx}
\usepackage{booktabs} 
\usepackage{amsmath}   
\usepackage{siunitx}   
\usepackage[hidelinks]{hyperref}   
\usepackage{caption}   
\usepackage{float}  
\usepackage{microtype}

\title{SCENARIO G CIVIL AIRWAYS}
\author{Celine He}
\date{March 2026}

\begin{document}

\maketitle

\tableofcontents

\section{Introduction}

\section{Vehicles}

\subsection{Vehicle Payloads}

\subsubsection{Drones}

\subsection{Initial sizing method (symbolic formulation)}


\begin{table}[h!]
\centering
\small
\renewcommand{\arraystretch}{1.2}
\begin{tabular}{ll}
\hline
Symbol & Meaning \\
\hline
$m_0,\ g,\ W$ & MTOW mass guess; gravity; weight $W=m_0 g$ \\
$\rho$ & Air density \\
$V_{\mathrm{stall}},\ V_{\mathrm{cruise}}$ & Stall speed; cruise speed \\
$S,\ (W/S)_\star,\ C_{L,\max}$ & Wing area; chosen wing loading; max lift coefficient \\
$A,\ (W/A)_\star,\ N,\ D$ & Total disk area; chosen disk loading; rotor count; rotor diameter \\
$C_{D0},\ k,\ C_L,\ C_D$ & Drag polar terms; lift and drag coefficients in cruise \\
$\eta_{\mathrm{p}},\ \eta_{\mathrm{pt}}$ & Propulsive efficiency; powertrain efficiency \\
$k_{\mathrm{int}},\ k_{\mathrm{pro}}$ & Hover loss factors (interference; profile) \\
$P_{\mathrm{h,elec}},\ P_{\mathrm{c,elec}}$ & Electrical hover power; electrical cruise power \\
$\lambda,\ \mathcal{E}_{\mathrm{tot}}$ & Reserve factor; total mission energy \\
\hline
\end{tabular}
\caption{Condensed symbol list for the sizing process.}
\end{table}

The initial sizing procedure begins with an assumed maximum take-off mass $m_0$ and combines two principal design parameters, namely the wing loading $W/S$ and the disk loading $W/A$, with the key mission requirements represented by the stall speed $V_\mathrm{stall}$ and cruise speed $V_\mathrm{cruise}$. This formulation is intended to generate first-order estimates of the wing area $S$, the total rotor disk area $A$, the rotor diameter $D$, and the electrical power demand in both hover and cruise.

\subsubsection{Weight definition}
\begin{equation}
W = m_0 g ,
\end{equation}
where $g$ is gravitational acceleration.

\subsubsection{Wing area estimation}
Two lower bounds for the wing reference area $S$ are evaluated in order to ensure compliance with both aerodynamic and design-loading requirements.

\paragraph{(a) Stall constraint}
At stall, the lift equals weight with $C_L=C_{L,\max}$:
\begin{equation}
W=\frac{1}{2}\rho V_\mathrm{stall}^2 S C_{L,\max}.
\end{equation}
Hence,
\begin{equation}
S_\mathrm{stall}=\frac{2W}{\rho V_\mathrm{stall}^2 C_{L,\max}} .
\end{equation}

\paragraph{(b) Wing-loading choice}
Given a selected wing loading $(W/S)_\star$,
\begin{equation}
S_{W/S}=\frac{W}{(W/S)_\star}.
\end{equation}

\paragraph{Governing wing area}
The sizing wing area is taken as
\begin{equation}
S=\max\left(S_\mathrm{stall},\, S_{W/S}\right).
\end{equation}

\subsubsection{Rotor disk area and diameter}
Given a selected disk loading $(W/A)_\star$, the total disk area is
\begin{equation}
A=\frac{W}{(W/A)_\star}.
\end{equation}
For $N$ identical rotors,
\begin{equation}
A_r=\frac{A}{N},
\end{equation}
where $A_r$ is the disk area per rotor. Assuming a circular disk $A_r=\pi D^2/4$, the rotor diameter becomes
\begin{equation}
D=\sqrt{\frac{4A_r}{\pi}}
=
\sqrt{\frac{4}{\pi}\frac{A}{N}}.
\end{equation}

\subsubsection{Hover power estimation}
Using actuator-disk momentum theory in hover, the ideal induced power is
\begin{equation}
P_{\mathrm{h,ideal}}=\frac{W^{3/2}}{\sqrt{2\rho A}}.
\end{equation}
To represent non-ideal effects (installation/interference losses and profile losses), multiplicative correction factors are applied:
\begin{equation}
P_{\mathrm{h}} = P_{\mathrm{h,ideal}}\,k_\mathrm{int}\,k_\mathrm{pro},
\end{equation}
where $k_\mathrm{int}>1$ and $k_\mathrm{pro}>1$.

Finally, converting to electrical power through the powertrain efficiency $\eta_\mathrm{pt}$,
\begin{equation}
P_{\mathrm{h,elec}}=\frac{P_{\mathrm{h}}}{\eta_\mathrm{pt}}.
\end{equation}

\subsubsection{Cruise power estimation}
At cruise speed $V_\mathrm{cruise}$, the dynamic pressure is
\begin{equation}
q=\frac{1}{2}\rho V_\mathrm{cruise}^2.
\end{equation}
The lift coefficient required for steady level cruise is
\begin{equation}
C_L=\frac{W}{qS}.
\end{equation}
Using a parabolic drag polar,
\begin{equation}
C_D = C_{D0}+kC_L^2,
\end{equation}
where $C_{D0}$ is the zero-lift drag coefficient and $k$ is the induced-drag factor.

The aerodynamic drag is
\begin{equation}
D_a=qSC_D,
\end{equation}
and the corresponding shaft power required is
\begin{equation}
P_{\mathrm{c,shaft}}=D_a V_\mathrm{cruise}.
\end{equation}
Converting to electrical power using propulsive efficiency $\eta_\mathrm{p}$ and powertrain efficiency $\eta_\mathrm{pt}$ gives
\begin{equation}
P_{\mathrm{c,elec}}=\frac{P_{\mathrm{c,shaft}}}{\eta_\mathrm{p}\eta_\mathrm{pt}}.
\end{equation}

\subsubsection{Outputs of initial sizing}
The initial sizing step therefore returns the set
\begin{equation}
\left\{S,\;A,\;D,\;P_{\mathrm{h,elec}},\;P_{\mathrm{c,elec}}\right\},
\end{equation}
which provides the first-cut geometry (wing area and rotor diameter) and the power levels that subsequently drive mission energy and battery/motor mass estimates in the overall MTOW iteration.






\bibliography{}

\end{document}
